% 第 1 篇：集合之美（编号与 ../lesson/chapters/01.tex 对应）
\corresponding{教案第 1 课《集合的含义与表示》}

\section{一个「自然」的问题}

19 世纪末，康托尔（Georg Cantor）创立了集合论。数学家一度以为，
终于可以为全部数学找到一个统一而严格的基础。然而没过多久，
一个看似简单的问题就让这座大厦动摇起来。

\begin{definition}[罗素悖论]
设
\[
  R=\{x\mid x\notin x\},
\]
即 $R$ 是所有「不以自身为元素」的集合组成的集合。
\end{definition}

\begin{problem}
$R$ 是不是 $R$ 的元素？
\end{problem}

\begin{solution}
分两种情况讨论：
\begin{itemize}
  \item 若 $R\in R$，则由 $R$ 的定义，$R$ 必须满足条件 $R\notin R$，矛盾；
  \item 若 $R\notin R$，则 $R$ 恰好满足定义中的条件，于是应有 $R\in R$，同样矛盾。
\end{itemize}
两种假设都导致矛盾。问题出在前提：我们默认了「所有集合的集合」
是一个合法的集合，而它并不是。
\end{solution}

这个悖论说明，「任意性质都能确定一个集合」这一直觉是错误的。
要修复它，必须对「什么样的集合可以存在」作出限制。

\section{公理化集合论：ZFC}

策梅洛（Zermelo）与弗兰克尔（Fraenkel）提出了一套公理，
把集合的存在性限制在若干条明确规则之内，这就是今天通用的
\key{ZFC 公理系统}。其中与中学内容最相关的几条是：

\begin{description}
  \item[外延公理] 两个集合相等，当且仅当它们含有相同的元素。
  \[
    A=B\iff(\forall x)(x\in A\leftrightarrow x\in B).
  \]
  \item[分离公理] 已知集合 $A$ 与性质 $p(x)$，可以取出
  $\{x\in A\mid p(x)\}$，它是一个集合。注意这里必须从已有的 $A$ 中「分离」，
  不能凭空构造「所有满足 $p(x)$ 的对象」。
  \item[幂集公理] 对任意集合 $A$，其所有子集组成的集合 $\mathcal{P}(A)$ 存在。
  \item[无穷公理] 存在一个无穷集合。这条公理保证了自然数集 $\mathbb{N}$ 的存在。
  \item[选择公理] 对任意一族两两不交的非空集合，可以从每个集合中
  各取一个元素组成新的集合。它看似显然，却与许多深刻结论相关。
\end{description}

\begin{example}
有了分离公理，罗素悖论就消失了：我们不能构造「所有满足 $x\notin x$
的对象」，只能从某个已有集合 $A$ 中分离出 $\{x\in A\mid x\notin x\}$。
而这个集合是否属于 $A$，是可以有确定答案的。
\end{example}

\section{无穷也有大小：对角线方法}

自然数集 $\mathbb{N}$ 是无穷的，实数集 $\mathbb{R}$ 也是无穷的。
它们一样大吗？康托尔的回答是：不一样。

\begin{definition}[等势]
若存在从集合 $A$ 到集合 $B$ 的一一对应，则称 $A$ 与 $B$ \key{等势}，
记作 $|A|=|B|$。
\end{definition}

\begin{definition}[可数集]
与自然数集 $\mathbb{N}$ 等势的集合称为 \key{可数集}。
\end{definition}

\begin{example}
整数集 $\mathbb{Z}$ 是可数的。按下表排列即可建立一一对应：
\[
  \begin{array}{c|ccccc}
  \mathbb{N} & 0 & 1 & 2 & 3 & \cdots\\ \hline
  \mathbb{Z} & 0 & 1 & -1 & 2 & \cdots
  \end{array}
\]
更意外的是，有理数集 $\mathbb{Q}$ 也是可数的——把分数按「分子分母之和」
分组排列，就能不重不漏地数完所有有理数。
\end{example}

那么实数呢？康托尔在 1874 年给出了否定的答案，用的正是著名的
\key{对角线方法}。

\begin{theorem}[康托尔]
开区间 $(0,1)$ 中的实数组成的集合不可数。
\end{theorem}

\begin{proof}
用反证法。假设 $(0,1)$ 可数，则可以把它列成一张无穷列表：
\[
  \begin{aligned}
  x_1&=0.a_{11}a_{12}a_{13}\ldots\\
  x_2&=0.a_{21}a_{22}a_{23}\ldots\\
  x_3&=0.a_{31}a_{32}a_{33}\ldots\\
  &\ \vdots
  \end{aligned}
  \qquad a_{ij}\in\{0,1,\dots,9\}.
\]
现在构造一个数 $y=0.b_1b_2b_3\ldots$，其中
\[
  b_k=
  \begin{cases}
    5, & a_{kk}\neq 5,\\
    6, & a_{kk}=5.
  \end{cases}
\]
这样 $y$ 与列表中的第 $k$ 个数在第 $k$ 位小数上不同，因此 $y$ 不等于
其中任何一个——但 $y\in(0,1)$，与「列表包含所有实数」矛盾。
故 $(0,1)$ 不可数。
\end{proof}

\begin{figure}[htbp]
  \centering
  \begin{tikzpicture}[scale=0.72]
    % 对角线示意
    \def\rows{5}
    \foreach \i in {1,...,\rows} {
      \foreach \j in {1,...,\rows} {
        \node at (\j*1.35, -\i*0.78) {$a_{\i\j}$};
      }
    }
    % 高亮对角线
    \foreach \i in {1,...,\rows} {
      \fill[second!25] ($(\i*1.35-0.55,-\i*0.78-0.32)$) rectangle ($(\i*1.35+0.55,-\i*0.78+0.32)$);
      \node[second] at (\i*1.35, -\i*0.78) {$a_{\i\i}$};
    }
    \draw[->, thick, second] (0.8, 0.4) -- (7.4, -3.6) node[pos=0.5, above, sloped] {对角线};
    \node at (8.6, -1.2) {$\longrightarrow\ b_k\neq a_{kk}$};
  \end{tikzpicture}
  \caption{对角线方法：沿对角线逐位改数字，得到的数不在原列表中。}
\end{figure}

对角线法的威力远不止于此。图灵用类似的思想证明了停机问题不可判定，
哥德尔用它构造了「本语句不可证明」的自指命题。它几乎是整个
数理逻辑中最重要的技巧之一。

\section{连续统假设：一个悬而未决的问题}

既然可数无穷与实数无穷不同，一个自然的问题是：
在这两者之间，还存在别的大小吗？

\begin{conjecture}[连续统假设]
不存在基数严格介于可数集与实数集之间的集合，即
\[
  2^{\aleph_0}=\aleph_1.
\]
\end{conjecture}

这个问题被列为希尔伯特 1900 年提出的 23 个问题的第一个。
它的结局出人意料：

\begin{itemize}
  \item 哥德尔在 1940 年证明，连续统假设与 ZFC 公理系统\key{相容}——
  在 ZFC 中无法证明它为假 \cite{godel1940}；
  \item 科恩在 1963 年证明，它的否定也与 ZFC \key{相容}——
  在 ZFC 中同样无法证明它为真 \cite{cohen1963}。
\end{itemize}

换言之，连续统假设在现有的集合论公理下\key{不可判定}。
要决定它，必须引入新的公理——而什么样的公理才算「正确」，
至今仍是一个开放的哲学与数学问题。

\section{结语}

从罗素悖论到 ZFC，从对角线法到连续统假设，集合论的历史几乎就是
一部「重新理解无穷」的历史。它告诉我们，直觉并不总是可靠，
而严格的公理化正是数学可靠性的来源。

\begin{quote}
「数学的本质在于它的自由。」——康托尔
\end{quote}
